\documentclass[11pt]{article}
\usepackage[margin=0.8in]{geometry}
\usepackage{enumitem}
\usepackage{setspace}
\setlist[itemize]{leftmargin=1.15em, itemsep=0.1em, topsep=0.18em}
\setlist[enumerate]{leftmargin=1.35em, itemsep=0.1em, topsep=0.18em}
\setlength{\parindent}{0pt}
\setlength{\parskip}{0.26em}
\setstretch{1.06}
\pagestyle{plain}

\begin{document}

\begin{center}
{\Large\bfseries Microeconomic Theory I}\\[0.3em]
Sample syllabus\\[0.15em]
Ke Cheng HSU\\
Department of Economics, National Chengchi University\\
Master's core, 3 credits, first semester, taught in English
\end{center}

\textbf{Place in the curriculum.}
This is the first half of the master's sequence Microeconomic Theory (I). Students who took sophomore Microeconomics have seen consumer theory, producer theory, and competitive markets. This course asks them to solve those problems with constrained optimization and to derive three duality results. General equilibrium, market power, and games are taught in Microeconomic Theory II. Mas-Colell, Whinston, and Green, \emph{Microeconomic Theory}, is reserved for the department's advanced graduate course and is not used here.

\textbf{Prerequisites.}
Intermediate undergraduate microeconomics and calculus, including partial derivatives and constrained optimization. The first problem set is a short review of the Lagrange method.

\textbf{Textbooks.}
Varian, Hal R., and Marc Melitz. \emph{Intermediate Microeconomics: A Modern Approach}. 10th edition. W.\ W.\ Norton, 2024. Calculus is required on the problem sets. The course reorders the book. Producer theory, Chapters 19--24, is taught before partial equilibrium, Chapters 15--16. The book places Chapter 16 before Chapter 19. Chapters 10--13 and 17--18 are omitted.

The expenditure function, Shephard's lemma, and Hotelling's lemma are taken from Varian, \emph{Microeconomic Analysis}, 3rd edition, Norton, 1992, and treated as envelope results. Those three results are required. The rest of that book is not. Nicholson and Snyder, \emph{Microeconomic Theory: Basic Principles and Extensions}, and Bergstrom and Varian, \emph{Workouts in Intermediate Microeconomics}, are sources of extra exercises. Lecture numbers follow Varian and Melitz.

\textbf{How the course runs.}
Each week opens with the optimization problem written on the board before the interpretation. Six problem sets are due at the start of class. A set asks students to set up the Lagrangian, solve for the demand or the supply, and sign one comparative static. Students may work in groups of up to three. Each student hands in a solution they can reproduce. On the day a set is due, two students present one question each. Solutions are posted the next day. Office hours are for a derivation that stopped partway, and students should bring the page where it stopped.

\textbf{Learning outcomes.}
By the end of the semester, students should be able to
\begin{itemize}
\item derive Marshallian and Hicksian demands and apply the Slutsky equation, including the sign of the own-price effect for a normal good,
\item state the expenditure function and Shephard's lemma, and recover a Hicksian demand from an expenditure function,
\item solve a cost-minimization problem and a profit-maximization problem, and state Hotelling's lemma,
\item find a partial-equilibrium competitive price and separate the efficiency claim in that market from the claim in a general-equilibrium economy.
\end{itemize}

\textbf{Assessment.}
Problem sets 25 percent. The presentation of a solution is part of that grade. Midterm examination 30 percent, covering weeks 1--7. Final examination 45 percent, cumulative for this course. Exams are closed book.

\textbf{Weekly plan.}
\begin{enumerate}
\item The market as a preview, and the budget set. Chapters 1--2. The numeraire, a kinked budget from a quantity discount, and what a tax does to the budget line.
\item Preferences. Chapter 3. Completeness, reflexivity, and transitivity. Monotonicity and convexity. Perfect substitutes, perfect complements, and Cobb--Douglas indifference curves.
\item Utility and choice. Chapters 4--5. A utility function as a label for the indifference map. The Lagrange condition: the marginal rate of substitution equals the price ratio. A corner solution.
\item Marshallian demand. Chapter 6. Homogeneity of degree zero in prices and income. Ordinary goods and Giffen goods. Quasilinear preferences and a zero income effect.
\item Revealed preference. Chapter 7. The weak axiom. Recovering a preference from a finite set of choices. Where revealed preference stops needing a utility function.
\item The Slutsky equation and Hicksian demand. Chapter 8. Substitution and income effects for a price change. The expenditure function as the least income that reaches a utility target. Shephard's lemma, from \emph{Microeconomic Analysis}: Hicksian demand is the derivative of the expenditure function.
\item Buying and selling. Chapter 9. An endowment budget. The income effect of a price change when the household sells the good. A labor-supply choice, leisure against consumption.
\item Midterm examination on weeks 1--7.
\item Technology and profit maximization. Chapters 19--20. Returns to scale. The profit function. Price equals marginal cost when the solution is interior. Hotelling's lemma: supply is the derivative of the profit function in the output price.
\item Cost minimization. Chapter 21. Conditional input demands. The cost function as the value function of that problem. Returns to scale read off the cost function.
\item Cost curves. Chapter 22. Fixed and sunk cost. Short-run and long-run average cost. Why the marginal-cost curve cuts the average-cost curve at its minimum.
\item Firm supply. Chapter 23. The short-run supply curve as marginal cost above average variable cost. Profit and producer surplus.
\item Industry supply. Chapter 24. Horizontal summation in the short run. Free entry and the long-run price in a constant-cost industry. Economic rent when firms differ.
\item Market demand. Chapter 15. Adding up individual demands. A representative consumer, and when that shortcut fails. Elasticity of market demand. The book places this chapter before Chapter 19. The course teaches it here, after firm supply.
\item Partial equilibrium. Chapter 16. Inverse demand and inverse supply. A specific tax, a quota, and a price ceiling, with the deadweight-loss area. Who bears the tax when supply is inelastic.
\item Consumer's surplus. Chapter 14. The integral under Marshallian demand as an approximation. Exact compensating and equivalent variation from the expenditure function, in one example. What a partial-equilibrium efficiency claim leaves out when other markets adjust.
\item Review. One consumer problem and one firm problem, each carried from the Lagrangian to the comparative static.
\item Final examination on this course.
\end{enumerate}

\end{document}
